\documentclass{article}
\usepackage[UTF8]{ctex}
\usepackage{amsmath,amssymb,amsfonts}
\usepackage{enumitem}
\usepackage{geometry}
\geometry{a4paper, margin=1in}
% === 错题本排版优化（不参与题目标识） ===
\setlength{\parindent}{0pt}
\setlength{\parskip}{0.35em plus 0.10em minus 0.08em}
\setlength{\emergencystretch}{3em}
\setlist{itemsep=0.25em, topsep=0.40em, parsep=0pt}
\allowdisplaybreaks[2]
% === 错题本排版优化结束 ===
\begin{document}
\section*{厦门大学2026年高等代数真题}
\section*{107．厦门大学2026年高等代数试卷}
一．填空题．每题5分，共50分．\\

\subsection{题目1}
1．设 \(n\) 阶方阵 \(A\) 的列分块为 \(A=\left(\alpha_{1}, \alpha_{2}, \cdots, \alpha_{n}\right), \operatorname{det}(A)=1\) ，令

\[
B=\left(\alpha_{1}-\alpha_{2}, 2 \alpha_{2}-\alpha_{3}, \cdots,(n-1) \alpha_{n-1}-\alpha_{n}, n \alpha_{n}-\alpha_{1}\right) .
\]

求 \(\operatorname{det}(B)=\)

\subsection{题目2}
2．设 4 阶矩阵 \(A=\left(\begin{array}{cccc}1 & -2 & 0 & 0 \\ 1 & 1 & 0 & 0 \\ 0 & 0 & 1 & 0 \\ 0 & 0 & 3 & 2\end{array}\right)\) ，其逆矩阵 \(A^{-1}=\) \(\_\_\_\_\)

\subsection{题目3}
3．向量 \(\alpha_{1}=(a, 1,-1,1), \alpha_{2}=(1,1, b, a), \alpha_{3}=(1, a, 1,-1)\) 。当 \(a=\) \(\_\_\_\_\)时，对任意 \(b\) 都使得向量组 \(\alpha_{1}, \alpha_{2}, \alpha_{3}\) 的秩为2。

\subsection{题目4}
4．设 \(n\) 阶方阵 \(A\) 满足 \(R(A)=4, R\left(A^{3}\right)=1\) ，且 \(R\left(A^{2}\right)>R\left(A^{3}\right)\) ，则 \(R\left(A^{2}\right)=\) \(\_\_\_\_\)．

\subsection{题目5}
5．设 \(\varphi\) 是 10 维线性空间 \(V\) 到 12 维线性空间 \(U\) 的线性映射，则 \(\operatorname{dim} \operatorname{Im} \varphi+\operatorname{dim} \operatorname{Ker} \varphi=\) \(\_\_\_\_\)．\\
\(\_\_\_\_\)

\subsection{题目6}
6．设多项式 \(f(x)=x^{5}+x^{4}-x^{3}+2 x^{2}-x-2\) ，则它在有理数域上的标准分解式为 \(\_\_\_\_\)．

\subsection{题目7}
7．设 3 阶方阵 \(A\) 有3个特征值 \(-1,-2,-3\) ，则 \(\operatorname{det}\left(A+A^{*}\right)=\) \(\_\_\_\_\)．

\subsection{题目8}
8．若 \(n\) 阶方阵 \(A\) 的元素全为 1，则其极小多项式为 \(\_\_\_\_\)．

\subsection{题目9}
9．设 \(A=\left(\begin{array}{lll}1 & 0 & 0 \\ 0 & 2 & 0 \\ 0 & 2 & 2\end{array}\right)\) ，则其 Jordan 标准形为 \(\_\_\_\_\)．\\
\(\_\_\_\_\)

\subsection{题目10}
10．设 \(f\left(x_{1}, x_{2}, x_{3}\right)=x_{1}^{2}+k x_{2}^{2}+8 x_{3}^{2}+2 x_{1} x_{2}+4 x_{1} x_{3}+4 k x_{2} x_{3}\) 是正定二次型，求 \(k\) 的取值范围\\
\(\_\_\_\_\)．

二．（15 分）设

\[
P=\left(\begin{array}{lll}
1 & 0 & 0 \\
1 & 1 & 0 \\
1 & 1 & 1
\end{array}\right), Q=\left(\begin{array}{lll}
0 & 1 & 1 \\
1 & 0 & 1 \\
1 & 1 & 0
\end{array}\right)
\]

满足 \(P X P+Q X Q=P X Q+Q X P+P^{2}-P Q\) ，求 \(X\) 。

三．（15分）设 \(A=\left(a_{i j}\right)_{3 \times 3}\) 为3阶实矩阵，\(a_{i j}=A_{i j}(i, j=1,2,3)\) ，若 \(a_{33}=1\) ，求线性方程组

\[
A\left(\begin{array}{l}
x_{1} \\
x_{2} \\
x_{3}
\end{array}\right)=\left(\begin{array}{l}
0 \\
0 \\
1
\end{array}\right)
\]

的解．

四．（15 分）设 \(A\) 为实矩阵满足 \(A^{2}=A\) ．证明：

\begin{itemize}
  \item[（1）] \(A\) 可对角化．
  \item[（2）] 存在实对称矩阵 \(B\) 和正定矩阵 \(C\) ，使得 \(A=B C\) 。
\end{itemize}

五。（15 分）（可能有误）设 \(f(x)\) 是实数域上的 \(n\) 次多项式 \((n \geq 3)\) ，若 \(\left(f(x), f^{\prime}(x)\right)=f^{\prime \prime}(x)\) ，且\(f(2025)=0\) ，求 \(f(x)\) ．

六。（15 分）设 \(f(x)=x^{n}+a_{1} x^{n-1}+\cdots+a_{n-1} x+a_{n}\) 是数域 \(F\) 上的不可约多项式，\(\varphi\) 是数域 \(F\) 上线性空间 \(V\) 上的线性变换，若 \(V\) 中非零向量 \(\alpha_{1}, \alpha_{2}, \cdots, \alpha_{n}\) 满足

\[
\varphi\left(\alpha_{i}\right)=\alpha_{i+1}(1 \leq i \leq n-1), \varphi\left(\alpha_{n}\right)=-a_{n} \alpha_{1}-a_{n-1} \alpha_{2}-\cdots-a_{1} \alpha_{n} .
\]

证明：\(\alpha_{1}, \alpha_{2}, \cdots, \alpha_{n}\) 线性无关。

七．（15 分）设 \(V\) 是复数域 \(\mathbb{C}\) 上的 \(n\) 维线性空间，\(\sigma\) 是 \(V\) 上的不可逆且非幂零的线性变换。证明：存在 \(\sigma-\) 子空间 \(V_{1}, V_{2}\) ，满足 \(V=V_{1} \oplus V_{2}\) ，且 \(\sigma \mid V_{1}\) 为可逆变换，\(\sigma \mid V_{2}\) 为幂零变换．

八．（10 分）证明：对任意的正整数 \(m, n\) ，必存在2025阶方阵 \(X\) ，使得

\[
X^{m}+X^{n}=\left(\begin{array}{cccccc}
2 & & & & & \\
2 & 2 & & & & \\
3 & 2 & 2 & & & \\
\vdots & \ddots & \ddots & \ddots & & \\
2024 & \ddots & 3 & 2 & 2 & \\
2025 & 2024 & \cdots & 3 & 2 & 2
\end{array}\right)
\]
\end{document}
